\section{System}
\label{sec:system}

We designed an agentic system to take text from a user in a conversational style, then use the resulting specifications to iteratively generate STLs until a passing STL is generated.

% ~1.5 pages. Describe what the app does today (the round's code). Keep history to one sentence per dropped idea.

\begin{figure*}[t]
  \centering
  \resizebox{\linewidth}{!}{% Figure 1: the pipeline. Pure TikZ (needs \usetikzlibrary{arrows.meta,positioning,calc}); included from
% sec/3_system.tex, and compiled on its own by figures/pipeline_preview.tex for checking.
\begin{tikzpicture}[
    font=\sffamily\footnotesize,
    >={Stealth[length=2mm]},
    box/.style={draw, rounded corners=2pt, align=center, minimum height=9mm, text width=26mm, inner sep=2pt},
    llm/.style={box, fill=blue!10, draw=blue!60!black},
    check/.style={box, fill=green!12, draw=green!45!black},
    ext/.style={box, fill=orange!15, draw=orange!70!black},
    plain/.style={box, fill=gray!8, draw=gray!60!black},
    io/.style={box, text width=15mm, fill=white, draw=black, thick},
    lab/.style={font=\sffamily\scriptsize, align=center, inner sep=1pt},
    flow/.style={->, thick},
    back/.style={->, thick, dashed},
  ]
  % --- row 1: build with code (CadQuery, or SDF for organic shapes) ---
  \node[io]    (req)  at (0, 0)     {Request\\(chat)};
  \node[llm]   (gate) at (3.0, 0)  {Intake gate\\scope, spec, missing sizes};
  \node[plain] (ret)  at (7.5, 0)  {Retrieve 3 examples\\CadQuery or SDF};
  \node[llm]   (gen)  at (10.6, 0)  {Write script\\+ expected size};
  \node[plain] (run)  at (13.7, 0) {Run in sandbox\\timeout, 2\,GB};
  \node[check] (chk)  at (16.8, 0) {Check mesh\\closed solid, size};
  \node[io]    (stl)  at (19.75, 0)  {STL};

  \draw[flow] (req) -- (gate);
  \draw[flow] (gate) -- node[lab, above] {in scope} node[lab, below] {or organic} (ret);
  \draw[flow] (ret) -- (gen);
  \draw[flow] (gen) -- (run);
  \draw[flow] (run) -- (chk);
  \draw[flow] (chk) -- node[lab, above] {pass} (stl);

  % repair loop (above row 1)
  \draw[back] (chk.north) -- ++(0, 0.55) -| node[lab, pos=0.25, above]
    {repair: error, failed check or measured size sent back (up to 4 attempts)} (gen.north);
  \draw[thick, dashed] (run.north) -- ++(0, 0.55);
  % clarify loop (above row 1, left)
  \draw[back] (gate.north) -- ++(0, 0.55) -| node[lab, pos=0.25, above]
    {clarify: ask the user} (req.north);

  % --- row 2: download ---
  \node[ext]   (ts)   at (7.5, -1.9)  {Thingiverse search\\gate's phrases};
  \node[check] (tn)   at (10.6, -1.9)  {Name / tag check\\is it the object?};
  \node[ext]   (td)   at (13.7, -1.9) {Download\\largest STL};

  \draw[flow] (gate.south) ++(-0.35, 0) |- node[lab, pos=0.8, above] {out of scope}
    node[lab, pos=0.8, below] {or sculpting failed} (ts.west);
  \draw[flow] (ts) -- (tn);
  \draw[flow] (tn) -- (td);
  \draw[flow] (td.east) -| (chk.south);

  % --- row 3: combined parts ---
  \node[plain] (plan) at (7.5, -3.8)  {Pieces plan\\from the gate};
  \node[plain] (each) at (10.6, -3.8)  {Build or download\\each piece (rows 1--2)};
  \node[plain] (pl)   at (13.7, -3.8) {Place by relation\\(on, beside, inside)};
  \node[plain] (mg)   at (16.8, -3.8) {Merge\\into one solid};

  \draw[flow] (gate.south) ++(0.35, 0) |- node[lab, pos=0.8, above] {composite} (plan.west);
  \draw[flow] (plan) -- (each);
  \draw[flow] (each) -- (pl);
  \draw[flow] (pl) -- (mg);
  \draw[flow] (mg.east) -| (stl.south);

  % --- legend ---
  \begin{scope}[shift={(0, -5.15)}, every node/.style={font=\sffamily\scriptsize, anchor=west}]
    \node[draw=blue!60!black, fill=blue!10, minimum size=2.5mm, inner sep=0, rounded corners=1pt] (l1) at (0, 0) {};
    \node at (0.25, 0) {LLM request};
    \node[draw=green!45!black, fill=green!12, minimum size=2.5mm, inner sep=0, rounded corners=1pt] at (2.3, 0) {};
    \node at (2.55, 0) {fixed check};
    \node[draw=orange!70!black, fill=orange!15, minimum size=2.5mm, inner sep=0, rounded corners=1pt] at (4.5, 0) {};
    \node at (4.75, 0) {outside source};
    \node[draw=gray!60!black, fill=gray!8, minimum size=2.5mm, inner sep=0, rounded corners=1pt] at (7.0, 0) {};
    \node at (7.25, 0) {code, no LLM};
    \draw[back] (9.4, 0) -- (10.2, 0);
    \node at (10.3, 0) {sent back for another try};
  \end{scope}
\end{tikzpicture}
}
  \caption{Pipeline. The intake gate (one LLM request) reads the chat and either asks for missing sizes or routes
    the request. In-scope parts are written as CadQuery code with three retrieved examples; organic objects are
    sculpted with SDF code through the same loop. Each script runs in a sandbox and the mesh is checked (closed
    solid, size against the request); a failure is sent back for a repair, up to four attempts. Objects code can't
    build are searched on Thingiverse, and combined parts are built or downloaded piece by piece, placed and merged.}
  \label{fig:pipeline}
\end{figure*}

\subsection{Intake gate}
% One LLM request: scope, filled-in spec, unknown parameters (critical or not). Fixed rules choose proceed /
% clarify / out of scope; composite plans with pieces. Not evaluated in this round (say so; Sec. 6).

A gate is required to determine the scope of the agentic development.
In this system, we define three different generation paths the agent can take: mechanical, organic, or composite.
A mechanical part requires exact dimensionality, whereas organic is more fluid in its design.
A composite request will combine multiple organic or mechanical requests into a single STL at the end.
It is up to the gate model to determine the scope of the user request, and to make clarifications until a valid request can be made.

The initial request by the user, $r_0$, is given to the \textit{Gate Model}.
From this initial request, the agent checks the request for any missing information, like height data or shape.
After, the agent drafts a request to be sent along in the pipeline, which the user can then suggest edits or approve.

\subsection{Retrieved examples}
% 26 hand-checked CadQuery examples (scripts/make_examples.py), each request plus generic shape keywords;
% TF-IDF retriever, top 3 into the prompt. History, one sentence: an embedding retriever, a similarity cutoff and
% an external Text-to-CadQuery pool were tried and dropped (no gain; retrieval-lib26.md, external-rag.md).

The agent then gathers similar examples to use as a basis for generating the STL.
There are 26 hand-crafted examples, ranging from gears to animals like fish.
Examples from Text-to-CadQuery were added initially as more examples, but more examples did not appear to improve performance.
Thus, those examples were dropped.

TF-IDF is used as the similarity metric, after initial testing showed embeddings to not improve overall performance.
TF-IDF is based on exact word matching, which was shown to be adequate for getting relevant examples.

All of the examples are scored, and the top 3 are appended to the request to allow the builder to see well-formed examples.

\subsection{Generate, run, check, repair}
% The model writes a script plus its own BBOX line (expected size). Sandbox: separate process, timeout, 2 GB cap.
% Checks: valid watertight mesh (trimesh), size matches the request. On failure the error goes back
% (script error, mesh check failed, size off with the measured size, #23; no code; reply cut off), up to 4 attempts.

The \textit{Builder Model} then gets the resulting output from the \textit{Gate Model}.
This model is instructed to build valid CadQuery~\cite{cadquery} (Python) code for the requested object.
The model also creates a suggested bounding box for the mesh, which is automatically checked against the actual bounding box of the resulting object.

The python code is ran in a sandbox that includes all the necessary libraries for performing generation.
If the code fails, the sandbox sends the error message and forces a rerun of the generation.
If it executes, it takes the resulting mesh and checks the sizing against the expected bounding box.
If those do not match within a percentage threshold, then the agent will also have to try again.

For organic parts, if configured, the model can instead make requests to an online service for parts.
In this case, Thingiverse~\cite{thingiverse} was configured as the online web service. It hosts user-made STLs that can be freely downloaded and used.
So, the model makes a search on its API for similar objects.
It can either pull the first match, or pull multiple and have the user decide.
If a Thingiverse model is selected, the task is considered complete.

\subsection{Combined parts}
% Composite plans: base first, pieces generated or downloaded, scaled to height, placed (on_top / under / beside /
% inside) by bounding boxes, merged with manifold3d. Not benchmarked; one example figure at most.

Composite, the third track, first decides what each part requires (mechanical or organic / downloaded).
It then runs each of those parts separately in the above desicribed generation loop.
After all the parts are gathered, it then composites the STLs into a final shape, merging them with manifold3d~\cite{manifold3d}.


\subsection{Model selection and fine-tuning}
\label{sec:finetune}
% Ollama for fully local use (default gemma4:12b). Student: Qwen3.5-4B, 4-bit QLoRA (rank 16, 2 epochs) on 1,052
% teacher conversations (822 CadQuery, 230 SDF; 990 train / 62 validation; teachers Sonnet and
% gemini-3.5-flash-lite), kept only when the part was valid
% and the right size (SDF: judge pass too). Failed attempts and repairs kept in the conversations. Requests are
% filtered against benchmark cases. Source: finetune-v2.md (data/training only; results from report-v1.md).

The agent is designed to work with multiple LLM orchestrators as the backend.
Primarily, API driven LLMs were used the most, but it was desirable to create a local LLM that the application could work with.

For the local LLM models, we took the newer Gemma 4 models and Qwen 3.5 and ran those as \textit{Gate Models} and \textit{Builder Models}.
We also wanted to see the effectiveness of fine-tuning a model specific to this task.
Fine tuning, like running these models, was performed locally. Thus, the options available were limited by hardware specifications.

The chosen model for fine-tuning was Qwen 3.5. The fine-tuning was performed by training a 4-bit QLoRA of rank 16 for 2 epochs.
It was trained on a generated 1052 teacher conversations: 822 CadQuery, 230 SDF; 990 train / 62 validation; teachers Sonnet and gemini-3.5-flash-lite.